\clearpage
\setlength{\headwidth}{\textwidth}
\addcontentsline{toc}{part}{PRIMERA PARTE: FONÉTICA}
\setcounter{chapter}{0}
\renewcommand{\thechapter}{\Roman{chapter}}
\setcounter{subsection}{19}
\setcounter{footnote}{0}
\refstepcounter{chapter}
\addcontentsline{toc}{chapter}{\thechapter\ -- Transcripción Fonética}
\begin{center}
  \underline{PRIMERA PARTE: FONÉTICA}\\[8.5mm]
  \underline{\thechapter\ -- Transcripción Fonética}
\end{center}
\vspace{4mm}

\subsection{} Los signos fonéticos empleados en este estudio son los de la
\emph{Revista de Filología Española}. Sin embargo, prescindimos, en las vocales,
de los matices abierto y cerrado, salvo allí donde éstos son significativos
para el habla de nuestra zona. Por lo demás, hay que hacer las aclaraciones
siguientes:

\begin{center}
\begin{tabular}{c@{${}={}$}l}
  \apt & a palatal \\
  \adp & a muy palatal, cerca de e abierta \\
  \textipa{\H{5}} & a relajada palatal \\
  \textipa{\`{@}} & a relajada muy palatal
\end{tabular}
\end{center}

(En el caso de las vocales \textipa{[e i o u]}, los signos diacríticos
\textipa{\textsubring{e}} o \textipa{\textsubwedge{e}} indican matiz abierto o
muy abierto.)

Dos puntos, encima de las vocales, indican matiz mixto y engolado
(cf. § 21).

A lo largo de este trabajo se indica la procedencia de cada forma por medio de
las siglas V (= Vega de Pas), P (= San Pedro del Romeral), S (= Selaya) o
R (= San Roque de Riomiera), seguidas de un número correspondiente al sujeto
informador (para estos números, cf. § 16). Cuando a la letra no le sigue
número, es que la forma ha sido proporcionada por una persona otra que los
informadores principales. Queremos subrayar que la localización de una forma
en uno, dos o tres puntos no excluye la posibilidad de que exista también en
otros puntos o en todos.

\clearpage
\refstepcounter{chapter}
\addcontentsline{toc}{chapter}{\thechapter\ -- Fonética Descriptiva}
\begin{center}
  \underline{\thechapter\ -- Fonética Descriptiva}
\end{center}
\vspace{4mm}
\section{\emph{Las vocales.}}

\subsection{} En la distribución entre las distintas sílabas del esfuerzo
espiratorio, el habla pasiega se une a las hablas del oeste peninsular en
general. El acento cae fuertemente en la sílaba tónica y por compensación las
sílabas átonas se relajan notablemente (cf. § 25).

Para la metafonía vocálica, cf. § 40 \emph{et seq.}

Las vocales que se transcriben \textipa{[a e i o u]} se pronuncian en general
como sus correspondientes castellanas. Sin embargo, habrá que considerar las
distintas series de sonidos que no tienen equivalencia en el idioma nacional.

\subsubsection*{i) vocales mixtas}

\setcounter{subsection}{20}
\subsection{ bis} \ic\ tónica y átona, se articula más interior y ligeramente
más abierta que \textipa{[i]}\footnote{El punto de articulación es casi
exactamente el de la \textipa{[i]} del inglés \emph{eat, sleep}.} y está
acompañada de cierto arrondamiento de los labios; es de timbre
engolado\footnote{Tónica, es el resultado de \textipa{[e]} o de
\textipa{[i]} inflexionadas (cf. §§ 42 y 44); inicial o intertónica, es el
resultado de \textipa{[e]} o de \textipa{[i]} átonas y se da en palabras que
tienen \ic, \uc\ o \apt\ tónicas (cf. § 25).}.

\uc\ tónica y átona, se pronuncia (paralelamente a \ic) adelantando algo el
punto de articulación y estrechando sensiblemente los labios; también es de
timbre engolado\footnote{Tónica, es el resultado de \textipa{[o]} o de
\textipa{[u]} metafonizadas (cf. §§ 43 y 45); inicial o intertónica, es el
resultado de \textipa{[o]} o de \textipa{[u]} átonas y se encuentra en
palabras que tienen \ic, \uc\ o \apt\ tónicas (cf. § 25).}.

\apt\ tónica y átona, se trata de una \apt\ palatal que se acerca en su punto
de articulación a una \emph{e} abierta, pero que se pronuncia ligeramente más
interior y más abierta que esta vocal y que se distingue de ella además por su
timbre engolado\footnote{La \ic\ y la \uc\ son sonidos bastante parecidos y
dan lugar fácilmente a sustituciones de una por otra, vgr.:
\dialectform{indruscu}{indr\rfeStressedMixedU sku} (V3) contra
\dialectform{andruisqui}{andrw\rfeStressedMixedI ski} (S1),
\dialectform{iscluiu}{iskl\rfeStressedMixedU j\rfeMixedU} (P2 R2) contra
\dialectform{iscluidu}{iskw\rfeStressedMixedI d\rfeMixedU} (V3 S1),
\dialectform{pidigüinu}{pidigw\rfeStressedMixedI n\rfeMixedU} `pedigüeño',
\dialectform{curpuñu}{k\rfeStressedMixedU rpun\rfeMixedU} `corpiño',
\dialectform{lubjisu}{lubj\rfeStressedMixedI s\rfeMixedU} `divieso'.}.

\subsubsection*{ii) vocales largas}

\subsection{} Éstas son \textipa{[\rfeStress{a}\rfeLong]} (vgr.:
\dialectform{preparái}{prepar\rfeStress{a}\rfeLong},
\dialectform{ixái}{iS\rfeStress{a}\rfeLong} `aguijón',
\dialectform{tostáis}{tost\rfeStress{a}\rfeLong s} `tostadas') e
\textipa{[\rfeStress{i}\rfeLong]}, que se oye
únicamente en la conjugación (vgr.:
\dialectform{merecís}{mereT\rfeStress{i}\rfeLong s} `merecéis',
\dialectform{agraicís}{agraiT\rfeStress{i}\rfeLong s} `agradecéis',
\dialectform{salís}{sal\rfeStress{i}\rfeLong s} `salís',
\dialectform{midís}{mid\rfeStress{i}\rfeLong s} `medís';
\dialectform{gulbí}{gulb\rfeStress{i}\rfeLong} `volved',
\dialectform{cumí}{kum\rfeStress{i}\rfeLong} `comed',
\dialectform{biní}{bin\rfeStress{i}\rfeLong} `venid';
\dialectform{cayístis}{kaj\rfeStress{i}\rfeLong st\rfeRelaxedClosed s}
`caísteis',
\dialectform{icístis}{iT\rfeStress{i}\rfeLong st\rfeRelaxedClosed s}
`dijisteis',
\dialectform{istubístis}{istub\rfeStress{i}\rfeLong st\rfeRelaxedClosed s}
`estuvisteis'). A veces se oye una \textipa{[u\rfeLong]} en posición final
(vgr.: \dialectform{mostru}{m\rfeStress{o}stru\rfeLong} `monstruo').

\subsubsection*{iii) las vocales finales}

\subsection{} Hay cuatro vocales finales, \textipa{[-\H{5}]},
\textipa{[-\'{5}]}, \textipa{[-u]} y \uc. Todas se articulan bastante más
relajadas que sus correspondientes castellanas.

\textipa{[-\rfeRelaxedPalatal]} es relajada y palatal; en Selaya y en San
Roque de Riomiera, suena muy palatal: \textipa{[-\rfeVeryPalatal]}. En la conjugación, precedida
de \textipa{[i]} acentuada, muchas veces suena \textipa{[-@]} (cf. § 133).
Vgr.: \dialectform{care}{k\rfeStress{a}r\rfeRelaxedPalatal},
\dialectform{mirjende}{mirj\rfeStress{e}nd\rfeRelaxedPalatal},
\dialectform{güene}{gw\rfeStress{e}n\rfeRelaxedPalatal} `buena',
\dialectform{negre}{n\rfeStress{e}gr\rfeRelaxedPalatal},
\dialectform{able}{\rfeStress{a}bl\rfeRelaxedPalatal} `habla',
\dialectform{ponge}{p\rfeStress{o}Ng\rfeRelaxedPalatal}. En posición
final no absoluta, conserva su matiz palatal, vgr.:
\dialectform{güenes}{gw\rfeStress{e}n\rfeRelaxedPalatal s},
\dialectform{negres}{n\rfeStress{e}gr\rfeRelaxedPalatal s},
\dialectform{ables}{\rfeStress{a}bl\rfeRelaxedPalatal s} `hablas',
\dialectform{pongen}{p\rfeStress{o}Ng\rfeRelaxedPalatal n}.

\textipa{[-\rfeRelaxedClosed]} se articula relajada y muy cerrada, pero sin
llegar casi nunca a \textipa{[i]} abierta:
\dialectform{grandi}{gr\rfeStress{a}nd\rfeRelaxedClosed},
\dialectform{dali}{d\rfeStress{a}l\rfeRelaxedClosed} `guadaña',
\dialectform{re}{r\rfeStress{e}\rfeRelaxedClosed} `red',
\dialectform{tardi}{t\rfeStress{a}rd\rfeRelaxedClosed},
\dialectform{sali}{s\rfeStress{a}l\rfeRelaxedClosed} (vb.),
\dialectform{comi}{k\rfeStress{o}m\rfeRelaxedClosed} (vb.),
\dialectform{intrjegi}{intrj\rfeStress{e}g\rfeRelaxedClosed} `entregue',
\dialectform{bibi}{b\rfeStress{i}b\rfeRelaxedClosed} `bebe' (imperat.),
\dialectform{curri}{k\rfeStress{u}r\rfeRelaxedClosed} `corre'
(imperat.). Final no absoluta, se puede abrir un poco, pero no
significativamente\footnote{Es decir, se diferencian fonéticamente las dos
realizaciones, pero no hay oposición fonológica ni morfológica.}:
\dialectform{grandis}{gr\rfeStress{a}nd\rfeRelaxedClosed s},
\dialectform{dalis}{d\rfeStress{a}l\rfeRelaxedClosed s},
\dialectform{tardis}{t\rfeStress{a}rd\rfeRelaxedClosed s},
\dialectform{salis}{s\rfeStress{a}l\rfeRelaxedClosed s},
\dialectform{comis}{k\rfeStress{o}m\rfeRelaxedClosed s},
\dialectform{intrjegin}{intrj\rfeStress{e}g\rfeRelaxedClosed n}.

\textipa{[-u]} fluctúa en su punto articulatorio entre una \textipa{[u]}
abierta y una \textipa{[o]} cerrada. La pronunciación general es
\textipa{[-u]}: \dialectform{malus}{m\rfeStress{a}lus},
\dialectform{manus}{m\rfeStress{a}nus},
\dialectform{cuernus}{kw\rfeStress{e}rnus},
\dialectform{estu}{\rfeStress{e}stu},
\dialectform{esto}{\rfeStress{e}sto},
(\dialectform{isti quesu es}{\rfeStress{i}sti k\rfeStress{e}su
\rfeStress{e}s}) \dialectform{güenu}{gw\rfeStress{e}nu},
\dialectform{güeno}{gw\rfeStress{e}no};
\dialectform{salgu}{s\rfeStress{a}lgu},
\dialectform{salgo}{s\rfeStress{a}lgo},
\dialectform{cantu}{k\rfeStress{a}ntu},
\dialectform{canto}{k\rfeStress{a}nto};
\dialectform{coxemus}{koS\rfeStress{e}mus},
\dialectform{salemus}{sal\rfeStress{e}mus} `salimos',
\dialectform{cantabamus}{kant\rfeStress{a}bamus},
\dialectform{quedremus}{kedr\rfeStress{e}mus} `querremos';
\dialectform{tubjendu}{tubj\rfeStress{e}ndu} `teniendo',
\dialectform{dixjendu}{diSj\rfeStress{e}ndu} `diciendo';
\dialectform{abrá biníu}{abr\rfeStress{a} bin\rfeStress{i}u}
`había venido'\footnote{En el § 40
se enumeran los casos donde aparece esta vocal final.}.

\uc\ se ha descrito en el § 21. Ejemplos son:
\dialectform{güinu}{gw\rfeStressedMixedI n\rfeMixedU} `bueno',
\dialectform{malu}{m\rfeStress{a}l\rfeMixedU},
\dialectform{agüilu}{agw\rfeStressedMixedI l\rfeMixedU} `abuelo',
\dialectform{surdu}{s\rfeStressedMixedU rd\rfeMixedU} `sordo',
\dialectform{sustu}{s\rfeStressedMixedU st\rfeMixedU},
\dialectform{punu}{p\rfeStressedMixedU n\rfeMixedU},
\dialectform{riyu}{r\rfeStressedMixedI j\rfeMixedU} `río' (sust.),
\dialectform{tiyu}{t\rfeStressedMixedI j\rfeMixedU}
`tío'\footnote{Por eso, esta vocal es morfema del masc. sing. (mientras el
neutro, cf. § 158, tiene la vocal final descrita inmediatamente antes).}.

\subsection{} Tras ciertos fonemas la final se ensordece parcial o totalmente.
Vgr. tras consonante sorda: \dialectform{estu}{\rfeStress{e}stu},
\dialectform{ece}{\rfeStress{e}T\rfeRelaxedPalatal} `hecha',
\dialectform{ace}{\rfeStress{a}T\rfeRelaxedPalatal} `haza',
\dialectform{puluque}{pul\rfeStress{u}k\rfeRelaxedPalatal},
\dialectform{isi}{\rfeStress{i}s\rfeRelaxedClosed} `ese',
\dialectform{pelebra}{pel\rfeStress{e}bra},
\dialectform{esu}{\rfeStress{e}su},
\dialectform{pocu}{p\rfeStress{o}ku},
\dialectform{axu}{\rfeStress{a}Su},
\dialectform{puzu}{p\rfeStress{u}Tu} `pozo',
\dialectform{pucu}{p\rfeStress{u}Tu} `podrido'; tras consonante sorda
precedida de otra sonora: \dialectform{tiempu}{tj\rfeStress{e}mbu},
\dialectform{cante}{k\rfeStress{a}nt\rfeRelaxedPalatal},
\dialectform{tranque}{tr\rfeStress{a}Nk\rfeRelaxedPalatal},
\dialectform{muertu}{mw\rfeStress{i}rtu},
\dialectform{suelte}{sw\rfeStress{e}lt\rfeRelaxedPalatal}; tras los grupos
\textipa{-tr-}, \textipa{-dr-}:
\dialectform{mayistru}{maj\rfeStress{i}stru} `maestro',
\dialectform{otre}{\rfeStress{o}tr\rfeRelaxedPalatal},
\dialectform{tulundru}{tul\rfeStress{u}ndru} `chichón'; tras
\textipa{[-r-]}: \dialectform{pere}{p\rfeStress{e}r\rfeRelaxedPalatal},
\dialectform{piru}{p\rfeStress{i}ru} `perro',
\dialectform{bure}{b\rfeStress{u}r\rfeRelaxedPalatal},
\dialectform{buru}{b\rfeStress{u}ru},
\dialectform{corre}{k\rfeStress{o}r\rfeRelaxedPalatal} `corre' (pres. ind.);
tras \textipa{[-n-]}:
\dialectform{rodanu}{rod\rfeStress{a}nu},
\dialectform{xontanu}{Sont\rfeStress{a}nu}.

Este fenómeno se produce igualmente cuando la vocal final va seguida de
\emph{-s}: \dialectform{estus}{\rfeStress{e}stus},
\dialectform{aces}{\rfeStress{a}T\rfeRelaxedPalatal s},
\dialectform{puluques}{pul\rfeStress{u}k\rfeRelaxedPalatal s},
\dialectform{pocus}{p\rfeStress{o}kus},
\dialectform{axus}{\rfeStress{a}Sus},
\dialectform{pucus}{p\rfeStress{u}Tus},
\dialectform{tiempus}{tj\rfeStress{e}mbus},
\dialectform{tranques}{tr\rfeStress{a}Nk\rfeRelaxedPalatal s},
\dialectform{muertus}{mw\rfeStress{i}rtus},
\dialectform{sueltis}{sw\rfeStress{e}lt\rfeRelaxedClosed s},
\dialectform{tulundrus}{tul\rfeStress{u}ndrus},
\dialectform{perus}{p\rfeStress{e}rus},
\dialectform{bures}{b\rfeStress{u}r\rfeRelaxedPalatal s},
\dialectform{corres}{k\rfeStress{o}r\rfeRelaxedPalatal s},
\dialectform{rodanus}{rod\rfeStress{a}nus}. En el caso de ser precedida y seguida de
\textipa{[-s-]}, por lo menos en la palabra aislada, la vocal puede perderse
totalmente, resultando una \textipa{[s:]} larga:
\dialectform{golós}{gol\rfeStress{o}s\rfeLong} `golosos',
\dialectform{ses}{s\rfeStress{e}s\rfeLong} `sesos',
\dialectform{iscuplós}{iskupl\rfeStress{o}s\rfeLong} `escrupulosos'.
